\documentclass[11pt]{article}
\usepackage[margin=1in]{geometry}
\usepackage{amsmath,amssymb,amsthm}
\usepackage{graphicx}
\usepackage[hidelinks]{hyperref}

\newtheorem{theorem}{Theorem}
\newtheorem{lemma}{Lemma}

\newcommand{\A}{\mathcal A}
\newcommand{\C}{\mathbb C}
\newcommand{\HH}{\mathcal H}
\newcommand{\K}{\mathcal K}
\newcommand{\J}{\mathcal J}
\newcommand{\N}{\mathbb N}
\newcommand{\Tr}{\operatorname{Tr}}
\newcommand{\ad}{\operatorname{ad}}

\title{A Norm-Closure Obstruction for Inner Jordan Derivations on \texorpdfstring{$K(H)$}{K(H)}}
\author{Run \texttt{fa\_banach\_001}, lane 14}
\date{2026-06-29}

\begin{document}
\maketitle

\noindent\textbf{Status.} \texttt{counterexample\_likely\_valid}.

\section{Source Problem}

The source paper is Jamjoom--Peralta--Siddiqui, \emph{Jordan weak
amenability and orthogonal forms on JB*-algebras}, arXiv:1411.2712
\cite{JamjoomPeraltaSiddiqui}. Remark 3.16 raises the problem of determining
whether, for a JB*-algebra $\J$, the inner Jordan derivations on $\J$ are
norm-dense in the set of all Jordan derivations on $\J$. The immediately
preceding Remark 3.15 recalls that for $\K(H)$, where $H$ is infinite
dimensional, there are derivations into $\K(H)^*$ which are not inner in the
Jordan sense.

\begin{center}
\includegraphics[width=0.97\linewidth]{figures/open_problem_crop.png}
\end{center}

\section{Proof Intuition}

For $\A=\K(\ell_2)$, every trace-class operator $T$ implements an associative
dual derivation
\[
  \delta_T(a)(b)=\Tr(T(ab-ba)).
\]
This is automatically a Jordan derivation when $\A$ is viewed with its Jordan
product. A finite inner Jordan derivation has a more special form: after
identifying $\A^*$ with the trace class, its implementer is a finite sum of
trace-class commutators, hence has trace zero.

The point is that the trace-zero obstruction survives norm closure in the
derivation norm. If $\delta_T$ is small, then $\Tr(T)$ is small. To see this,
move a large finite-dimensional block of $\ell_2$ to a disjoint far-out block
by a compact partial isometry. Evaluating $\delta_T$ on that partial isometry
and its adjoint compares the trace of $T$ on the two blocks; the first block
captures $\Tr(T)$ and the far-out block has negligible trace.

\section{Counterexample}

\begin{theorem}
Let $H=\ell_2(\N)$ and let $\A=\K(H)$, regarded as a JB*-algebra with Jordan
product $a\circ b=(ab+ba)/2$. Then the inner Jordan derivations from
$\A$ into $\A^*$ are not norm-dense in the set of all Jordan derivations from
$\A$ into $\A^*$.
\end{theorem}

\begin{proof}
Identify $\A^*$ isometrically with the trace class $S_1(H)$ by
$T\mapsto \psi_T$, where $\psi_T(x)=\Tr(Tx)$ for $x\in\K(H)$. For
$T\in S_1(H)$ define
\[
  \delta_T(a)(b)=\Tr(T(ab-ba)) \qquad (a,b\in\A).
\]
Equivalently, $\delta_T(a)$ is represented by the trace-class operator
$Ta-aT$. Hence
\[
  \|\delta_T\|=\sup_{\|a\|\leq 1}\|Ta-aT\|_1 .
\]
Since $\delta_T$ is an associative derivation from the $C^*$-algebra $\A$
into the dual bimodule $\A^*$, it is a Jordan derivation for the Jordan
product:
\[
  \delta_T(a\circ b)=\delta_T(a)\circ b+a\circ\delta_T(b).
\]

We first record the form of the inner Jordan derivations. Let
$S\in S_1(H)$, let $c\in\K(H)$, and let $\phi=\psi_S$. The elementary inner
Jordan derivation determined by $\phi$ and $c$ is
\[
  D_{\phi,c}=[L(\phi),L(c)] : \A\to\A^*.
\]
For $a,b\in\A$, the computation in the source paper gives
\[
  D_{\phi,c}(a)(b)
  =\phi\bigl(b\circ(c\circ a)-(c\circ b)\circ a\bigr).
\]
Using $x\circ y=(xy+yx)/2$ in the associative algebra $\K(H)$,
\[
  b\circ(c\circ a)-(c\circ b)\circ a
  =\frac{1}{4}[c,[a,b]].
\]
Therefore
\[
  D_{\phi,c}(a)(b)
  =\frac{1}{4}\Tr(S[c,[a,b]])
  =\frac{1}{4}\Tr((Sc-cS)[a,b])
  =\delta_{(Sc-cS)/4}(a)(b),
\]
where we used cyclicity of the trace for products of trace-class and bounded
operators. The implementer $(Sc-cS)/4$ has trace zero. Hence every finite sum
of inner Jordan derivations is of the form $\delta_R$ with $\Tr(R)=0$.

Next we prove that trace is continuous with respect to the derivation norm on
these implemented derivations. Fix $T\in S_1(H)$ and an orthonormal basis
$(e_k)_{k\geq 1}$ of $H$. Let $p_n$ be the projection onto
$\operatorname{span}\{e_1,\ldots,e_n\}$. Then $\Tr(Tp_n)\to\Tr(T)$. Since
$T$ is trace class, after passing far enough into the basis we can choose a
rank-$n$ projection $q_n$ orthogonal to $p_n$ such that $\Tr(Tq_n)\to 0$.
Let $v_n$ be a partial isometry with $v_n^*v_n=p_n$ and $v_nv_n^*=q_n$.
Then $v_n\in\K(H)$ and $\|v_n\|=\|v_n^*\|=1$. Thus
\[
  \delta_T(v_n)(v_n^*)
  =\Tr(T(v_nv_n^*-v_n^*v_n))
  =\Tr(T(q_n-p_n))\longrightarrow -\Tr(T).
\]
Consequently
\[
  |\Tr(T)|\leq \|\delta_T\|.
\]

Let $P$ be a rank-one projection. Then $\Tr(P)=1$, so $\delta_P$ is a Jordan
derivation with nonzero trace implementer. If $D$ is a finite inner Jordan
derivation, then $D=\delta_R$ for some $R\in S_1(H)$ with $\Tr(R)=0$, and the
previous estimate gives
\[
  \|\delta_P-D\|=\|\delta_{P-R}\|\geq |\Tr(P-R)|=1.
\]
The same lower bound holds for every norm limit of finite inner Jordan
derivations. Therefore $\delta_P$ is not in the norm closure of the inner
Jordan derivations from $\A$ into $\A^*$.
\end{proof}

\section{Verification Notes}

The proof is analytic and uses no computation. The only source-specific
calculation is the standard formula for a finite inner Jordan derivation,
already displayed in Remark 3.15 of the source paper; the packet expands it
in the special associative algebra $\K(H)$ and observes that its trace-class
implementer has trace zero.

Novelty checks were bounded as follows. The run indexes were searched for
\texttt{1411.2712}, \texttt{Jordan weak amenability}, \texttt{inner Jordan
derivations}, \texttt{norm dense}, \texttt{norm-dense}, \texttt{K(H)}, and
\texttt{trace}. Later local papers citing \cite{JamjoomPeraltaSiddiqui} were
checked around Jordan and triple derivation approximation statements. Web
searches on 2026-06-29 for exact and close norm-density variants found no
separate exact resolution.

\section{Limitations and Review Recommendation}

This counterexample answers the general norm-density problem negatively by
one concrete JB*-algebra, $\K(\ell_2)$. It does not classify which JB*-
algebras have norm-dense inner Jordan derivations, and it does not address
the separate normal/predual approximation questions for JBW*-algebras
mentioned at the end of Remark 3.16.

Human review should focus on the trace-continuity estimate
$|\Tr(T)|\leq \|\delta_T\|$ and on the identification of finite inner Jordan
derivations with trace-zero implementers.

\begin{thebibliography}{9}
\bibitem{JamjoomPeraltaSiddiqui}
F. B. Jamjoom, A. M. Peralta, and A. A. Siddiqui,
\emph{Jordan weak amenability and orthogonal forms on JB*-algebras},
arXiv:1411.2712, 2014.
\end{thebibliography}

\end{document}
